\chapter{Die Prüfung, die verschwindet}\label{fall:pruefung}

\begin{merksatz}
\boxtitel{Die Vorlage}
\textbf{Zusammenhang.} Messerfassung, C. Prüft, ob die Summe zweier Rohwerte im zulässigen Bereich liegt. Läuft $500$ mal je Sekunde.

\begin{center}
\begin{minipage}{0.9\textwidth}
\code{int mittelwert(int a, int b) \{}\\
\code{\ \ \ \ int s = a + b;}\\
\code{\ \ \ \ if (s < 0) \{}\\
\code{\ \ \ \ \ \ \ \ melde\_ueberlauf();}\\
\code{\ \ \ \ \ \ \ \ return FEHLER;}\\
\code{\ \ \ \ \}}\\
\code{\ \ \ \ return s / 2;}\\
\code{\}}
\end{minipage}
\end{center}

\begin{center}
\begin{minipage}{0.88\textwidth}
\kursiv{Acht Zeilen mit einer ausdrücklichen Überlaufprüfung. Der Verfasser hat an den Fall gedacht, ihn behandelt und eine Meldung vorgesehen --- nach allen Maßstäben sorgfältige Arbeit.}
\end{minipage}
\end{center}
\end{merksatz}

\begin{merksatz}
\boxtitel{Die Diagnose}
\begin{center}
\renewcommand{\arraystretch}{1.5}
\begin{tabular}{@{}P{1.2cm}P{7.4cm}@{}}
\toprule
$1$ & Überlauf bei \code{int} ist undefiniertes Verhalten.\\
$2$ & Der Übersetzer darf annehmen, dass er nicht vorkommt.\\
$3$ & Bei $a, b \ge 0$ und ohne Überlauf ist $s \ge 0$.\\
$4$ & \textbf{Also ist \code{s < 0} immer falsch.}\\
$5$ & \textbf{Die Prüfung wird entfernt.}\\
\bottomrule
\end{tabular}
\end{center}

\begin{center}
\begin{minipage}{0.88\textwidth}
\kursiv{Die Prüfung soll einen Überlauf abfangen und \kursiv{ist} selbst der Überlauf. Sie steht im Quelltext und nicht im übersetzten Programm --- Band 21, Kapitel~7. Der Übersetzer hat nichts falsch gemacht.}
\end{minipage}
\end{center}

\textbf{Der Nachweis.}

\begin{center}
\renewcommand{\arraystretch}{1.45}
\begin{tabular}{@{}P{2.2cm}P{2.0cm}P{2.4cm}P{2.2cm}@{}}
\toprule
\textbf{Stufe} & \textbf{Befehle} & \textbf{Prüfung da?} & \\
\midrule
\code{-O0} & $24$ & \textcolor{gruen}{\textbf{ja}} & \textbf{im Versuch}\\
\code{-O1} & $9$ & \textcolor{gruen}{ja} & \\
\code{-O2} & $\mathbf{4}$ & \textcolor{akzent}{\textbf{nein}} & \textcolor{akzent}{\textbf{Aus\-lieferung}}\\
UBSan & --- & \textcolor{gruen}{meldet} & \textcolor{gruen}{$0{,}4$ s Prüflauf}\\
\bottomrule
\end{tabular}
\end{center}

\begin{center}
\kursiv{Der Code wurde mit \code{-O0} geprüft und mit \code{-O2} ausgeliefert. Alle Prüfungen waren grün, und im ausgelieferten Programm fehlte die Sicherung --- der Fehler war zwei Jahre im Bestand.}
\end{center}
\end{merksatz}

\begin{technik}
\boxtitel{Das Urteil: umgebaut}
\begin{center}
\begin{minipage}{0.9\textwidth}
\code{int mittelwert(int32\_t a, int32\_t b, int32\_t *aus) \{}\\
\code{\ \ \ \ int64\_t s = (int64\_t)a + (int64\_t)b;\ \ /* kein Ueberlauf */}\\
\code{\ \ \ \ if (s < MESS\_MIN * 2 || s > MESS\_MAX * 2)}\\
\code{\ \ \ \ \ \ \ \ return -ERANGE;}\\
\code{\ \ \ \ *aus = (int32\_t)(s / 2);}\\
\code{\ \ \ \ return 0;}\\
\code{\}}
\end{minipage}
\end{center}

\begin{center}
\begin{minipage}{0.86\textwidth}
\kursiv{Drei Änderungen: breiterer Rechentyp, damit der Überlauf nicht eintreten \kursiv{kann}; Prüfung gegen fachliche Grenzen statt gegen das Vorzeichen; Ergebnis über Zeiger, Fehler als Rückgabewert. Der Umbau betrifft die Signatur --- deshalb umgebaut und nicht überarbeitet.}
\end{minipage}
\end{center}
\end{technik}

\begin{merksatz}
\boxtitel{Die Begründung}
\begin{center}
{\small
\renewcommand{\arraystretch}{1.45}
\begin{tabular}{@{}P{2.8cm}P{2.0cm}P{2.2cm}P{2.0cm}@{}}
\toprule
& \textbf{Vorlage} & \textbf{Umgebaut} & \\
\midrule
Zeilen & $8$ & $7$ & \\
Undefiniertes Verhalten & \textcolor{akzent}{\textbf{ja}} & \textcolor{gruen}{\textbf{nein}} & \\
Prüfung überlebt \code{-O2} & \textcolor{akzent}{nein} & \textcolor{gruen}{ja} & \\
Prüft gegen & \textcolor{akzent}{Vorzeichen} & \textcolor{gruen}{\textbf{Fachgrenzen}} & \\
Aufrufstellen geändert & --- & $\mathbf{4}$ & $0{,}5$ PT\\
Laufzeit & $0{,}4$ ns & $0{,}6$ ns & \textcolor{gruen}{unerheblich}\\
\bottomrule
\end{tabular}}
\end{center}

\begin{center}
\begin{minipage}{0.88\textwidth}
\kursiv{Die vierte Zeile ist der eigentliche Gewinn und geht über die Behebung hinaus: Die Vorlage prüfte gegen das Vorzeichen --- eine technische Eigenschaft. Der Umbau prüft gegen den Messbereich des Geräts, und das fängt auch Werte ab, die nicht überlaufen und trotzdem unmöglich sind.}
\end{minipage}
\end{center}

\textbf{Was daraus für den Bestand folgte.}

\begin{center}
\renewcommand{\arraystretch}{1.45}
\begin{tabular}{@{}P{3.0cm}P{2.2cm}P{2.6cm}@{}}
\toprule
\textbf{Maßnahme} & \textbf{Aufwand} & \textbf{Funde}\\
\midrule
Diese Funktion & $0{,}5$ PT & $1$\\
\textbf{UBSan über den Prüfbestand} & \textcolor{gruen}{$\mathbf{0{,}5}$ \textbf{PT}} & \textcolor{akzent}{$\mathbf{6}$ \textbf{weitere}}\\
Prüfstrecke mit \code{-O2} & $1$ PT & \textcolor{gruen}{dauerhaft}\\
\bottomrule
\end{tabular}
\end{center}

\begin{center}
\kursiv{Ein halber Personentag für den Sanitizer, sechs weitere Fälle derselben Art. Das ist die Rechnung aus Band 21, Kapitel~10 --- und der Grund, warum dieser Fall nicht mit der Behebung endet.}
\end{center}
\end{merksatz}

\begin{warnung}
\boxtitelr{Der Fehlschluss}
\begin{center}
\renewcommand{\arraystretch}{1.5}
\begin{tabular}{@{}P{3.2cm}P{2.2cm}P{2.4cm}@{}}
\toprule
\textbf{Naheliegend} & \textbf{Richtig?} & \\
\midrule
\enquote{Der Übersetzer ist schuld} & \textcolor{akzent}{\textbf{nein}} & \textbf{Band 21, Kap.~9}\\
\enquote{Nie optimieren} & \textcolor{akzent}{nein} & \textcolor{akzent}{Faktor $10$}\\
\enquote{Immer \code{int64\_t}} & \textcolor{akzent}{\textbf{nein}} & \textcolor{akzent}{doppelter Speicher}\\
\textbf{\enquote{Mit denselben Einstellungen prüfen wie ausliefern}} & \textcolor{gruen}{\textbf{ja}} & \\
\bottomrule
\end{tabular}
\end{center}

\begin{center}
\begin{minipage}{0.88\textwidth}
\kursiv{Die letzte Zeile ist die eigentliche Lehre und keine Sprachfrage: Ein Prüfbestand, der eine andere Übersetzung prüft als die ausgelieferte, prüft ein anderes Programm. Das gilt für Optimierungsstufen genauso wie für Konfigurationen und Bibliotheksversionen.}
\end{minipage}
\end{center}

\begin{center}
\kursiv{Die dritte Zeile ist die Übertreibung, die teuer wird: Bei einem Gerät mit vierundsechzig Kilobyte verdoppelt ein breiterer Typ den Speicherbedarf jedes Puffers. Der breitere Typ gehört in die \kursiv{Rechnung}, nicht in die Speicherung.}
\end{center}
\end{warnung}
\vspace{-4pt}
\noindent Lattner\zit{20} erklärt, wie der Op\-ti\-mie\-rer aus un\-de\-fi\-nier\-tem Ver\-hal\-ten schließt; Chisnall\zit{5} zeigt, wie weit das C-Modell von der Ma\-schi\-ne ent\-fernt ist.
